The EU eFTI (Electronic Freight Transport Information) regulation requires digital exchange of freight data
between competent authorities and transport operators. By 2027, Member States must be ready to accept
freight information in digital form and make it available to competent authorities.
As a result, each Member State must have its own eFTI gateway to transmit information between parties.

According to the regulation, eDelivery is the basis for data exchange between Member States and within national systems.
For cross-border exchange, static eDelivery configuration is used,
where required certificates and other parameters are configured manually.
The eFTI gateway must also provide an eDelivery interface to transport operators (eFTI platforms),
but it may also offer other integration options, such as a REST interface.

The eDelivery solution provided to platforms must use dynamic configuration, which automates
the discovery of certificates and other technical parameters.
For dynamic metadata discovery, eDelivery uses two additional components:
SMP (Service Metadata Publisher) and SML (Service Metadata Locator).

This thesis studies how eDelivery dynamic discovery works in practice and implements a working prototype.
The work analyzes existing open-source components, identifies practical integration limitations,
and implements a minimal Kotlin-based SMP/SML solution that supports Harmony Access Point
dynamic discovery requirements.

The resulting prototype demonstrates full dynamic message exchange between two access points and is evaluated
with functional tests and performance measurements. The results show that a focused implementation with minimal
dependencies can provide high throughput for both SMP metadata responses and SML DNS lookups in a local test setup.


% No need to change this
The thesis is in \langEng~and contains \calculatepages pages of text, 
\total{totalchapters} chapters\ifthenelse{\equal{\totvalue{figure}}{0}}{}{% If no figures, do nothing
, \total{figure} \ifnum\totvalue{figure}=1 figure\else figures\fi%
}\ifthenelse{\equal{\totvalue{table}}{0}}{}{% If no tables, do nothing
, \total{table} \ifnum\totvalue{table}=1 table\else tables\fi%
}.
